\begin{thebibliography}{99}
\bibitem{SmithTate}
S. P. Smith and J. Tate,
\emph{The center of the 3-dimensional and 4-dimensional Sklyanin algebras},
K-Theory \textbf{8} (1994), 19--63.
中心三次元の表示はp.~35, Remark (1)。
\url{https://doi.org/10.1007/BF00962089}.
\bibitem{Walton}
C. Walton, \emph{Representation theory of three-dimensional Sklyanin algebras},
Nuclear Physics B \textbf{860} (2012), 167--185.
Theorem 2.14、式(2.15)--(2.16)。著者公開版p.~9を確認。
\url{https://math.rice.edu/~notlaw/RepThySkly_NucPhys_revised_final.pdf}.
\bibitem{DeLaetLeBruyn}
K. De Laet and L. Le Bruyn,
\emph{The geometry of representations of 3-dimensional Sklyanin algebras},
Algebras and Representation Theory \textbf{18} (2015), 761--776.
確認した版はarXiv:1405.1158v1、Section 2, p.~3。
中心式は同箇所の対称な再表示を用いる。
\url{https://arxiv.org/abs/1405.1158}.
\bibitem{Orlov}
D. Orlov, \emph{Derived categories of coherent sheaves and triangulated categories of singularities},
arXiv:math/0503632v3 (2009), Theorem 2.5。
非可換なnoetherian連結次数付きGorenstein代数に適用する。
\url{https://arxiv.org/abs/math/0503632}.
\bibitem{CCKM}
T. Cassidy, A. Conner, E. Kirkman and W. F. Moore,
\emph{Periodic free resolutions from twisted matrix factorizations},
Journal of Algebra \textbf{455} (2016), 137--163.
Proposition 5.7、Theorem 5.8およびfinite projective resolutionによる比較。
\url{https://doi.org/10.1016/j.jalgebra.2016.01.037};
\url{https://users.wfu.edu/moorewf/Papers/twistedMFs.pdf}.
\bibitem{EG}
P. Etingof and V. Ginzburg,
\emph{Noncommutative del Pezzo surfaces and Calabi--Yau algebras},
Journal of the European Mathematical Society \textbf{12} (2010), 1371--1416.
Theorem 3.6.2、Sections 3.5--3.6。genericな仮定を無条件に拡大して用いない。
\url{https://arxiv.org/abs/0709.3593}.
\bibitem{Elagin}
A. Elagin, \emph{On equivariant triangulated categories},
arXiv:1403.7027v2 (2015), Theorem 4.2およびSection 8。
有限アーベル群の双対同変化とdg enhancementを用いる。
\url{https://arxiv.org/abs/1403.7027}.
\bibitem{LZ}
X. Lin and S. Zhang,
\emph{Serre algebra, matrix factorization and categorical Torelli theorem for hypersurfaces},
arXiv:2310.09927v2 (2024), Theorem 1.7, p.~3。
本稿では可換模型の一意性にのみ使用する。
\url{https://arxiv.org/abs/2310.09927}.
\bibitem{ArtebaniDolgachev}
M. Artebani and I. Dolgachev, \emph{The Hesse pencil of plane cubic curves},
arXiv:math/0611590v3 (2008), Lemma 1。
\url{https://arxiv.org/abs/math/0611590}.
\bibitem{BurkeStevenson}
J. Burke and G. Stevenson, \emph{The derived category of a graded Gorenstein ring},
arXiv:1507.00830v1 (2015), Definition 5.1、Theorem 6.4。
正則な底の相対Orlov比較を検討するための文献であり、一般量子化族への全適用条件は未確認。
\url{https://arxiv.org/abs/1507.00830}.
\end{thebibliography}